% answers_short.tex —— 2003 年·春季上海卷 · 简短答案速查
% 由 tools/make_short_answers.py 自 answers.tex 的 \gksolhead/\gkanswer 提取，
% 并与 2003/2003_数学真题_春季上海_答案要点.md 逐题比对后生成，勿手改（改源 MD 后重跑）。

\begin{gknotice}
\textbf{使用说明}：本节只给最终答案便于快速核对，完整推导见解析版。答案为编写者独立推导，\textbf{非官方标准答案}。
\end{gknotice}

\gkpart{填空题}
\begin{itemize}[leftmargin=2.6em]
  \item[\textbf{1.}] $4$。
  \item[\textbf{2.}] $\dfrac{\pi}{3}$（$60^\circ$）。
  \item[\textbf{3.}] 第二象限。
  \item[\textbf{4.}] $(3,2)$。
  \item[\textbf{5.}] $a\le -2$。
  \item[\textbf{6.}] $\operatorname{Re}z>1$。
  \item[\textbf{7.}] $a=4\pm\sqrt5$。
  \item[\textbf{8.}] $x\in(0,\dfrac{\pi}{2})$。
  \item[\textbf{9.}] $16$ 场。
  \item[\textbf{10.}] $\arctan\dfrac{3}{8}$。
  \item[\textbf{11.}] 非退化时 $b=6$；若允许退化区间，另有 $a=b=1$。
  \item[\textbf{12.}] $3\sqrt2$。
\end{itemize}

\gkpart{选择题}
\begin{center}\begin{tabular}{ccccc}
\toprule
题号 & 13 & 14 & 15 & 16 \\
\midrule
答案 & D & A & C & A \\
\bottomrule
\end{tabular}\end{center}

\gkpart{解答题}
\begin{itemize}[leftmargin=2.6em]
  \item[\textbf{17.}] $(1,2)\cup(4,5)$。
  \item[\textbf{18.}] $\left(\dfrac{\pi}{6}+4k\pi,\sqrt3\right)$ 或 $\left(\dfrac{5\pi}{6}+4k\pi,\sqrt3\right)$，$k\in\mathbb Z$。
  \item[\textbf{19.}] （1）正三棱柱；（2）$45^\circ$。
  \item[\textbf{20.}] （1）奇函数；在 $(-\infty,0)$、$(0,+\infty)$ 上递增；（2）两式均为 $0$，恒等式 $f(x^2)=5f(x)g(x)$。
  \item[\textbf{21.}] （1）$\dfrac{x^2}{4}+\dfrac{y^2}{3}=1$；（2）$\dfrac{(2X+1)^2}{4}+\dfrac{4Y^2}{3}=1$；（3）$k_{PM}k_{PN}=\dfrac{b^2}{a^2}$。
  \item[\textbf{22.}] （1）A：$1500+230(n-1)$元，B：$2000(1.05)^{n-1}$元；（2）选A；（3）最大约$827$元，发生在第19年。
\end{itemize}
